\documentclass[11pt]{article}
\usepackage[utf8]{inputenc}
%\usepackage[english]{babel}
\usepackage[OT1]{fontenc}
\usepackage[margin=0.5in]{geometry}
\usepackage{lmodern}
\usepackage{amssymb}
\usepackage{amsmath}
\usepackage{amsfonts}
\usepackage{eurosym}
\usepackage{setspace}
% \doublespacing
\usepackage{comment}
\usepackage{caption}
\usepackage{pdflscape}
\usepackage{subfigure}
\usepackage{array}
\usepackage{amsthm}
\usepackage[bottom]{footmisc}
\usepackage{hyperref}
\usepackage{float}
\usepackage{listingsutf8}
\usepackage{graphicx}
\usepackage[dvipsnames]{xcolor}
\usepackage{stackengine}
\newcommand\overlinebelow[1]{\stackunder[3pt]{$#1$}{\rule{.9ex}{.055ex}}}

%\makeatletter
%\renewcommand{\fnum@figure}{Gráfico \thefigure}
%\makeatother
%\makeatletter
%\renewcommand{\fnum@table}{Tabla \thetable}
\makeatother
\usepackage{cases}
 \setcounter{MaxMatrixCols}{20}


\usepackage{mathrsfs}  
\newcommand{\R}{\Rightarrow}  
\newcommand{\Z}{\mathbb{Z}}
\newcommand{\N}{\mathbb{N}}
\newcommand{\Q}{\mathbb{Q}}
\newcommand{\E}{\mathbb{E}}



\usepackage{enumitem}   % for customized enumerate/itemize labels








\begin{document}

\title{Trade Final Project}


\author{Yoon, Yuxuan, Tomas}


\date{Fall 2026}

\maketitle
\tableofcontents{}

\large

\newpage

\section{Background: Bianchi (2011)}

\textbf{Preferences:}
\[
\mathbb{E}_0\sum_{t=0}^{\infty}\beta^t u(c_t),
\]
where the consumption aggregator is
\[
c_t
=
\left[
\omega (c_t^T)^{-\eta}
+
(1-\omega)(c_t^N)^{-\eta}
\right]^{-\frac{1}{\eta}}.
\]

\textbf{Budget constraint:}
\[
b_{t+1}+c_t^T+p_t^N c_t^N
=
(1+r)b_t+y_t^T+p_t^N y_t^N.
\]

\textbf{Credit constraint:}
\[
b_{t+1}
\geq
-\left(
\kappa^N p_t^N y_t^N
+
\kappa^T y_t^T
\right).
\]

$$y_{i,t}^{T} =
    y_{c,t}^T + \varepsilon_{i,t}^T  \quad 
    i=\text{poor,rich}$$

\textbf{Household optimality conditions:}
\[
\lambda_t=u'(c_t)\frac{\partial c_t}{\partial c_t^T},
\]

\[
p_t^N
=
\frac{1-\omega}{\omega}
\left(
\frac{c_t^T}{c_t^N}
\right)^{\eta+1},
\]

\[
\lambda_t
=
\beta(1+r)\mathbb{E}_t[\lambda_{t+1}]
+
\mu_t,
\]

\[
b_{t+1}
+
\kappa^N p_t^N y_t^N
+
\kappa^T y_t^T
\geq 0,
\]

\[
\mu_t
\left[
b_{t+1}
+
\kappa^N p_t^N y_t^N
+
\kappa^T y_t^T
\right]
=0.
\]

\textbf{Market clearing:}
\[
c_t^N=y_t^N,
\]

\[
c_t^T
=
y_t^T+(1+r)b_t-b_{t+1}.
\]

\textbf{Equilibrium price of non-tradables:}
\[
p_t^N
=
\frac{1-\omega}{\omega}
\left(
\frac{c_t^T}{y_t^N}
\right)^{\eta+1}.
\]

\textbf{Recursive household problem:}
\[
V(b,B,y)
=
\max_{b',c^T,c^N}
\left\{
u(c(c^T,c^N))
+
\beta\mathbb{E}_{y'|y}
V(b',B',y')
\right\},
\]

subject to
\[
b'
+
p^N(B,y)c^N
+
c^T
=
y^T+b(1+r)+p^N(B,y)y^N,
\]

\[
b'
\geq
-\left[
\kappa^N p^N(B,y)y^N
+
\kappa^T y^T
\right],
\]

\[
B'=\Gamma(B,y).
\]

\textbf{Social planner's problem:}
\[
V(b,y)
=
\max_{b',c^T}
\left\{
u(c(c^T,y^N))
+
\beta\mathbb{E}_{y'|y}V(b',y')
\right\},
\]

subject to
\[
b'+c^T=y^T+b(1+r),
\]

\[
b'
\geq
-\left[
\kappa^N
\frac{1-\omega}{\omega}
\left(
\frac{c^T}{y^N}
\right)^{\eta+1}
y^N
+
\kappa^T y^T
\right].
\]

\textbf{Planner's first-order conditions:}
\[
\lambda_t^{sp}
=
u_t^T
+
\mu_t^{sp}\Psi_t,
\]

\[
\lambda_t^{sp}
=
\beta(1+r)
\mathbb{E}_t[\lambda_{t+1}^{sp}]
+
\mu_t^{sp},
\]

where
\[
\Psi_t
\equiv
\kappa^N
\frac{p_t^N c_t^N}{c_t^T}
(1+\eta).
\]

\textbf{Key externality:}
\[
\underbrace{
\beta(1+r)
\mathbb{E}_t
\left[
\mu_{t+1}^{sp}\Psi_{t+1}
\right]
}_{\text{external marginal cost of borrowing}}
\]

\[
\text{Private agents do not internalize }
\beta(1+r)
\mathbb{E}_t
\left[
\mu_{t+1}^{sp}\Psi_{t+1}
\right],
\]

\[
\boxed{
\text{Competitive equilibrium}
\;\Rightarrow\;
\text{overborrowing}
}
\]

\[
\boxed{
\text{Overborrowing}
=
\text{debt in the decentralized equilibrium}
-
\text{debt under the constrained-efficient allocation}.
}
\]

\newpage



\section{Setting}

\begin{itemize}
    \item \textbf{Agents:}
\[
i\in\{P,R\}
\equiv
\{\text{poor},\text{rich}\}.
\]

Let $\omega_i$ denote the population weight of household $i$, with $\omega_P+\omega_R=1$. For simplicity, we can set
\[
\omega_P=\omega_R=\frac{1}{2}.
\]

\item \textbf{Preferences:}

Each household $i\in\{P,R\}$ has preferences given by
\[
\mathbb{E}_0
\sum_{t=0}^{\infty}
\beta^t u(c_{i,t}), \qquad c_{i,t}
=
\left[
\omega (c_{i,t}^{T})^{-\eta}
+
(1-\omega)(c_{i,t}^{N})^{-\eta}
\right]^{-1/\eta}
\]
where the consumption basket is a CES aggregator of tradable and
nontradable goods, with $\eta>-1$, $\omega\in(0,1)$

\item \textbf{Endowments:}

Each household receives tradable and nontradable endowments.

Tradable income is given by
\[
y_{i,t}^{T}
=
y_{c,t}^{T}
+
\varepsilon_{i,t}^{T},
\qquad
i\in\{P,R\},
\]
where $y_{c,t}^{T}$ is the common, aggregate component of tradable income and
$\varepsilon_{i,t}^{T}$ is an idiosyncratic shock.

We assume that nontradable income is common across households:
\[
y_{i,t}^{N}=y_t^{N},
\qquad
i\in\{P,R\}.
\]

\item \textbf{Debt:}

Each household has access to a one-period, non-state-contingent bond
denominated in units of tradable goods.

$b_{i,t+1}$ denote household $i$'s bond holdings chosen at time $t$.

The aggregate bond holdings are
\[
B_t
=
\sum_{i\in\{P,R\}}
\omega_i b_{i,t}.
\]
With equal population weights,
\[
B_t
=
\frac{1}{2}b_{P,t}
+
\frac{1}{2}b_{R,t}.
\]

\item \textbf{Budget constraint:}

For each household $i$,
\[
b_{i,t+1}
+
c_{i,t}^{T}
+
p_t^N c_{i,t}^{N}
=
(1+r)b_{i,t}
+
y_{i,t}^{T}
+
p_t^N y_{i,t}^{N}.
\]

\item \textbf{Credit constraint:}

Each household faces a collateral constraint:
\[
b_{i,t+1}
\geq
-
\left(
\kappa^N p_t^N y_{i,t}^{N}
+
\kappa^T y_{i,t}^{T}
\right).
\]
\end{itemize}



\section{Household Problem}

HH $i$ problem:
\begin{align*}
    \max_{\left\{
c_{i,t}^{T},
c_{i,t}^{N},
b_{i,t+1}
\right\}_{t=0}^\infty} \quad & \mathbb{E}_0
\sum_{t=0}^{\infty}
\beta^t u(c_{i,t}) \\ 
\text{s.t.} \quad &  c_{i,t}
=
\left[
\omega (c_{i,t}^{T})^{-\eta}
+
(1-\omega)(c_{i,t}^{N})^{-\eta}
\right]^{-1/\eta}  \\
& b_{i,t+1} +c_{i,t}^{T}+p_t^N c_{i,t}^{N}=(1+r)b_{i,t}+y_{i,t}^{T}+p_t^N y_{i,t}^{N} \\ 
& b_{i,t+1}
\geq
-
\left(
\kappa^N p_t^N y_{i,t}^{N}
+
\kappa^T y_{i,t}^{T}
\right)
\end{align*}

Let $\lambda_{i,t}$ denote the multiplier on the budget constraint and
$\mu_{i,t}$ the multiplier on the borrowing constraint.

FOC:
\begin{align*}
    [c_{i,t}^T]: \quad &  \lambda_{i,t} = u_{T}^i(t) \ \ \R \ \ \lambda_{i,t}
=
u'(c_{i,t})
\frac{\partial c_{i,t}}{\partial c_{i,t}^{T}} = u'(c_{i,t}) \cdot (-1/\eta) \cdot \Big[ \cdot\Big]^{-1/\eta-1}\cdot \omega \cdot (-\eta) \cdot (c_{i,t}^{T})^{-\eta-1}
\\  \\ 
  [c_{i,t}^N]: \quad &  \lambda_{i,t}p_t^N = u_{N}^i(t) \ \ \R \ \ \lambda_{i,t}
=\frac{u'(c_{i,t})
\dfrac{\partial c_{i,t}}{\partial c_{i,t}^{N}}}{p_t^N} 
= \frac{u'(c_{i,t}) \cdot (-1/\eta) \cdot \Big[ \cdot\Big]^{-1/\eta-1}\cdot (1-\omega) \cdot (-\eta) \cdot (c_{i,t}^{N})^{-\eta-1}}{p_t^N}  \\  \\ 
[b_{i,t+1}]: \quad & \left\{\begin{array}{lcc}
\lambda_{i,t}=\beta(1+r)\mathbb{E}_t\left[
\lambda_{i,t+1}\right]+\mu_{i,t} \\  \\ 
b_{i,t+1} + \kappa^N p_t^N y_{i,t}^{N} +\kappa^T y_{i,t}^{T}\geq 0 \\  \\ 
\mu_{i,t}\geq 0 \\  \\ 
\mu_{i,t} \left[b_{i,t+1}+\kappa^N p_t^N y_{i,t}^{N}+\kappa^T y_{i,t}^{T}\right]=0    
\end{array}\right.
\end{align*}

Or, combining  $[c_{i,t}^T]$ with  $[c_{i,t}^N]$:
\begin{align*}
    [c_{i,t}^T,c_{i,t}^N]: \quad &  p_t^N
=
\frac{1-\omega}{\omega}
\left(
\frac{c_{i,t}^{T}}
{c_{i,t}^{N}}
\right)^{\eta+1}
\\  \\  
[b_{i,t+1}]: \quad & \left\{\begin{array}{lcc}
\lambda_{i,t}=\beta(1+r)\mathbb{E}_t\left[
\lambda_{i,t+1}\right]+\mu_{i,t} \\  \\ 
b_{i,t+1} + \Big(\kappa^N p_t^N y_{i,t}^{N} +\kappa^T y_{i,t}^{T}\Big)\geq 0, \quad =0 \ \ \text{if $\mu_{i,t}>0$}  
\end{array}\right.
\end{align*}

 

 

\section{Market clearing}

\begin{itemize}
    \item The nontradable goods market clears domestically:
\[
\sum_{i\in\{P,R\}}
\omega_i c_{i,t}^{N}
=
\sum_{i\in\{P,R\}}
\omega_i y_{i,t}^{N}.
\]

Since $y_{i,t}^{N}=y_t^N$:
\[
\sum_{i\in\{P,R\}}
\omega_i c_{i,t}^{N}
=
y_t^N.
\]

\item Tradable goods market condition:
\[
\sum_{i\in\{P,R\}}
\omega_i c_{i,t}^{T}
+
B_{t+1}
=
(1+r)B_t
+
Y_t^T,
\]
where 
\[
B_t \equiv \sum_{i\in\{P,R\}}\omega_i b_{i,t}.
\]
and, using $y_{i,t}^{T}=y_{c,t}^{T}+\varepsilon_{i,t}^{T}$:
\[
Y_t^T
\equiv 
\sum_{i\in\{P,R\}}
\omega_i y_{i,t}^{T} = y_{c,t}^{T}
+
\sum_{i\in\{P,R\}}
\omega_i\varepsilon_{i,t}^{T}
\]


\end{itemize}



\section{Equilibrium Definition - Decentralized RCE}


We consider the optimization problem of each household in recursive form. 

Aggregate bond holdings are given by $B=\omega_P b_P+\omega_R b_R$.

The aggregate endowment vector is
$$\mathbf{y} \equiv (y^N,y_c^T),$$
and the vectors of the distribution of idiosyncratic income shocks and of bond holdings are
$$\mathbf{e}\equiv (\varepsilon_P^T,\varepsilon_R^T) \quad , \quad \mathbf{b} \equiv (b_P,b_R) $$

Thus, the current aggregate state is
\[
\mathbf{s}=\left(\mathbf{y},\mathbf{e},\mathbf{b}\right),
\]
Unlike in the representative-agent version of Bianchi (2011), aggregate
bond holdings alone may not be sufficient to characterize the aggregate
state. In particular, the distribution of individual bond holdings and
idiosyncratic income shocks across households can affect aggregate
consumption, the demand for nontradables, and therefore the equilibrium
price of nontradables. Hence, the distribution of individual states is
also relevant for forecasting future aggregate outcomes, which is crucial
to form expectations of the price of nontradables.

We denote by $\Gamma(\cdot)$ the perceived law of motion for the distribution of bond
holdings. Given the current aggregate state $\mathbf{s}$, households perceive the distribution of bond holdings next period to be
\[
\mathbf{b}'=\Gamma(\mathbf{s}).
\]
The forecast price of nontradables is therefore given by a pricing function
\[
p^N=p^N(\mathbf{s}).
\]

The relevant individual state variables for household $i$ include its
individual bond holdings and its idiosyncratic income shock.

The recursive problem of household $i$ can be written as
\begin{align*}
V_i(b_i,\varepsilon_i^T;\mathbf{s})
=
\max_{b_i',\,c_i^T,\,c_i^N} \quad & u\left(c_i(c_i^T,c_i^N)\right) +\beta \, 
\mathbb{E}_{\mathbf{y}',\,\varepsilon_i^{T\prime}\mid\mathbf{y},\,\varepsilon_i^T}\left[V_i\left(b_i',\varepsilon_i^{T\prime};\mathbf{s}'\right)\right] \\ 
\text{s.t.} \quad &b_i'+p^N(\mathbf{s})c_i^N+c_i^T=y_i^T+b_i(1+r)+p^N(\mathbf{s})y^N, \quad \text{where \ } y_i^T=y_c^T+\varepsilon_i^T \\ 
& b_i'\geq-\left(\kappa^N p^N(\mathbf{s})y^N+\kappa^T y_i^T\right) \\ 
& \mathbf{b}'=\Gamma(\mathbf{s})
\end{align*}


The solution to the household problem yields decision rules for individual
bond holdings, tradable consumption and nontradable consumption
\[
\widehat b_i(b_i,\varepsilon_i^T;\mathbf{s}), \quad 
\widehat c_i^T(b_i,\varepsilon_i^T;\mathbf{s}), \quad 
\widehat c_i^N(b_i,\varepsilon_i^T;\mathbf{s}),
\qquad
i\in\{P,R\},
\]
with associated value functions
\[
V_i(b_i,\varepsilon_i^T;\mathbf{s}).
\]

The household optimization problem induces a mapping from the perceived
law of motion for the distribution of bond holdings to an actual law of motion,
given by the individual households' choices $\Big[\widehat b_P(b_P,\varepsilon_P^T;\mathbf{s}) \ \ , \ \
\widehat b_R(b_R,\varepsilon_R^T;\mathbf{s})\Big]$.
In a rational expectations equilibrium, the perceived and actual laws of
motion must coincide.


\ 

\ 

\ 

\textbf{Definition (Decentralized Recursive Competitive Equilibrium)}

A decentralized recursive competitive equilibrium for this two-agent SOE is defined by a pricing function
\[
p^N(\mathbf{s}),
\]
a perceived law of motion
\[
\Gamma(\mathbf{s}),
\]
and decision rules
\[
\Big\{ \ 
\widehat b_i(b_i,\varepsilon_i^T;\mathbf{s}) \ , \ 
\widehat c_i^T(b_i,\varepsilon_i^T;\mathbf{s}) \ , \ 
\widehat c_i^N(b_i,\varepsilon_i^T;\mathbf{s}) \ 
\Big\}_{i\in\{P,R\}},
\]
with associated value functions
\[
\left\{
V_i(b_i,\varepsilon_i^T;\mathbf{s})
\right\}_{i\in\{P,R\}},
\]
such that the following conditions hold:

\begin{enumerate}

\item[(i)] \textbf{Household optimization:}
For each $i\in\{P,R\}$,
\[
\left\{
\widehat b_i(b_i,\varepsilon_i^T;\mathbf{s}),
\widehat c_i^T(b_i,\varepsilon_i^T;\mathbf{s}),
\widehat c_i^N(b_i,\varepsilon_i^T;\mathbf{s}),
V_i(b_i,\varepsilon_i^T;\mathbf{s})
\right\}
\]
solve the recursive optimization problem of household $i$ for given
$p^N(\mathbf{s})$ and $\Gamma(\mathbf{s})$.

\item[(ii)] \textbf{Rational expectations condition:}
The perceived law of motion for the distribution of bond holdings is consistent
with the actual law of motion:
\[
\Gamma(\mathbf{s})
= \Big[ \ \widehat b_P(b_P,\varepsilon_P^T;\mathbf{s})\ , \ 
\widehat b_R(b_R,\varepsilon_R^T;\mathbf{s}) \ \Big]
\]


\item[(iii)] \textbf{Markets clear:}

\begin{align*}
    [MC.N]: \quad & \sum_{i\in\{P,R\}}\omega_i \cdot \widehat c_i^N(b_i,\varepsilon_i^T;\mathbf{s})=y^N \\
     [MC.T]: \quad &  \underbrace{\sum_{i\in\{P,R\}}\omega_i \cdot \widehat b_i^N(b_i,\varepsilon_i^T;\mathbf{s})}_{B'} +\sum_{i\in\{P,R\}}\omega_i \cdot \widehat c_i^T(b_i,\varepsilon_i^T;\mathbf{s})=B(1+r)+Y^T
\end{align*}
\end{enumerate}

 


\end{document}
